%! Graphing Exponential Functions and Transformations — Unit 6, Lesson 6.2 — Exit Ticket (Answer Key)
\documentclass[10pt]{article}
\usepackage{algebra2-article}
\usepackage{algebra2-key}   % pulls in -boxes; do NOT also load -boxes
\newcommand{\TallMath}[1]{$\displaystyle #1\rule[-1.4em]{0pt}{3.2em}$}
\newcommand{\lgrid}[4]{%
  \draw[step=1, linegray!40, very thin] (#1,#3) grid (#2,#4);
  \draw[->, semithick, charcoal] (#1,0)--(#2,0) node[below right, font=\scriptsize]{$x$};
  \draw[->, semithick, charcoal] (0,#3)--(0,#4) node[left, font=\scriptsize]{$y$};}
% #1=a  #2=ln(b)  #3=h  #4=k  #5=xmin  #6=xmax   (ln2 = 0.693147)
\newcommand{\expgraph}[6]{\draw[very thick, forest, domain=#5:#6, samples=140, smooth]
  plot (\x,{(#1)*exp((#2)*((\x)-(#3)))+(#4)});}
\newcommand{\hasym}[3]{\draw[navy, dashed, semithick] (#1,#3)--(#2,#3);}

\begin{document}

\pageheader{Unit 6, Lesson 6.2}{Exit Ticket --- Answer Key}
\namedateperiod

\vspace{0.08in}

No notes. For every function, find the asymptote \emph{before} you write the range.

\begin{enumerate}[leftmargin=1.8em, itemsep=9pt]
  \item \textbf{Read it off the equation.} \; $y=4^{x}-5$

        \vspace{3pt}
        Asymptote: $y=$ \ans{$-5$} \quad Range: \ans{$(-5,\infty)$} \quad
        $y$-intercept: \ans{$(0,-4)$} \quad Growth or decay? \ans{growth}

  \item \textbf{Inside or outside.} Only one of these has range $(3,\infty)$.

        \vspace{3pt}
        \hspace*{1.2em} $y=2^{\,x+3}$ \qquad\qquad $y=2^{x}+3$

        \vspace{3pt}
        Which one? \ans{$y=2^{x}+3$} \; Explain in one sentence what the \emph{other} one does instead.

        \vspace{2pt}
        \ansline{$y=2^{\,x+3}$ shifts the graph left $3$; the asymptote stays at $y=0$, so its range is still $(0,\infty)$.}

  \item \textbf{Write the equation} of the curve below. It has the form $y=a\,b^{x}+k$.
    \begin{center}
    \begin{minipage}[c]{0.44\linewidth}
    \centering
    \begin{tikzpicture}[scale=0.42]
      \lgrid{-3}{4}{-1}{10}
      \begin{scope}\clip (-3,-1) rectangle (4,10); \expgraph{2}{0.693147}{0}{1}{-3}{4} \end{scope}
      \hasym{-3}{4}{1}
      \node[navy, font=\scriptsize] at (3,1.8) {$y=1$};
      \fill[charcoal] (0,3) circle (3pt) node[below left, font=\scriptsize]{$(0,3)$};
      \fill[charcoal] (1,5) circle (3pt) node[right, font=\scriptsize]{$(1,5)$};
    \end{tikzpicture}
    \end{minipage}
    \begin{minipage}[c]{0.52\linewidth}
    $k=$ \ans{1} \quad $a=$ \ans{2} \quad $b=$ \ans{2} \\[8pt]
    $y=$ \ans{$2(2)^{x}+1$}
    \end{minipage}
    \end{center}

  \item \textbf{SOL-style multiple choice.} What is the range of $g(x)=3(2)^{x}-6$?
        \begin{enumerate}[label=\Alph*., leftmargin=1.9em, itemsep=1pt]
          \item $(-6,\infty)$ \quad \textcolor{keyred}{\textbf{$\leftarrow$ correct}}
          \item $[-6,\infty)$
          \item $(0,\infty)$
          \item $(6,\infty)$
        \end{enumerate}
\end{enumerate}

\begin{teachernote}
\textbf{Answer 4: A.} Each distractor is a named error. \textbf{B} is Lesson 6.0's bracket error --- the
asymptote is a boundary the range \emph{excludes}, so $-6$ is never attained. \textbf{C} ignored the
vertical shift entirely and reported the parent's range. \textbf{D} took the sign of $k$ from the
equation's ``$-6$'' backwards.

\textbf{Sort the collected tickets into four piles} to decide how Lesson 6.3 opens:
\emph{got it} / \emph{inside-outside} (item 2 answered $2^{\,x+3}$, or item 1's range given as
$(0,\infty)$) / \emph{bracket} (item 4 answered B, or item 1's range written $[-5,\infty)$) /
\emph{grabbed the intercept} (item 3 with $a=3$). The inside-outside pile is the only one that needs a
full re-teach; run it off the Hook's graphs II and III, where the same $3$ does two different things.

\textbf{Item 3 is the hardest item and the most informative.} A student who writes $a=3$ read the
$y$-intercept as $a$ and forgot that the intercept is $a+k$. That is a one-minute fix but it will not
fix itself --- reteach the order: \emph{dashed line first, then intercept minus $k$, then one more
point}. Lesson 6.5 leans on this same reading.

\textbf{Item 1's ``growth'' is a deliberate check} that students are still reading growth versus decay
off $b$ and not off the sign attached to the function. The base is $4$, so the function grows even
though the $y$-intercept is negative.
\end{teachernote}

\end{document}
